\documentclass[11pt]{article}

% ============================================================================
%  The Ising Model
%  Structured LaTeX transcription of Chapter 3 of
%    S. Friedli & Y. Velenik, "Statistical Mechanics of Lattice Systems:
%    a Concrete Mathematical Introduction" (Cambridge University Press, 2017).
%
%  This file reproduces the chapter's organization, numbered definitions and
%  theorem statements, and connecting exposition, for use as a reference while
%  formalizing the model in Lean. Proofs are summarized, not transcribed
%  verbatim.  Equation/result numbering follows the source.
% ============================================================================

\usepackage[utf8]{inputenc}
\usepackage[T1]{fontenc}
\usepackage{amsmath,amssymb,amsthm,mathtools,mathrsfs}
\usepackage[margin=1in]{geometry}
\usepackage{hyperref}

% --- theorem environments --------------------------------------------------
\theoremstyle{plain}
\newtheorem{theorem}{Theorem}
\renewcommand{\thetheorem}{3.\arabic{theorem}}
\newtheorem{lemma}[theorem]{Lemma}
\newtheorem{proposition}[theorem]{Proposition}
\newtheorem{corollary}[theorem]{Corollary}
\theoremstyle{definition}
\newtheorem{definition}[theorem]{Definition}
\newtheorem{exercise}[theorem]{Exercise}
\theoremstyle{remark}
\newtheorem{remark}[theorem]{Remark}

% --- notation --------------------------------------------------------------
\newcommand{\Z}{\mathbb{Z}}
\newcommand{\R}{\mathbb{R}}
\newcommand{\C}{\mathbb{C}}
\newcommand{\Rgeq}{\R_{\ge 0}}
\newcommand{\fin}{\Subset}          % finite subset of Z^d
\newcommand{\incr}{\Uparrow}        % van Hove / increasing convergence to Z^d
\newcommand{\La}{\Lambda}
\newcommand{\sg}{\sigma}
\newcommand{\Ham}{\mathcal{H}}
\newcommand{\Eset}{\mathscr{E}}
\newcommand{\bdy}{\#}
\newcommand{\avg}[1]{\langle #1 \rangle}
\newcommand{\meanp}[1]{\langle #1 \rangle^{+}}
\newcommand{\meanm}[1]{\langle #1 \rangle^{-}}
\newcommand{\bc}{\mathscr{B}}
% Sections render as 3.1--3.10 and results as 3.x, mirroring the book's
% Chapter 3; the theorem counter is advanced manually where the book's
% numbering has gaps (results not transcribed here).
\renewcommand{\thesection}{3.\arabic{section}}

\title{\textbf{The Ising Model}\\[4pt]
\large Structured reference transcription of Chapter 3 of\\
Friedli \& Velenik, \emph{Statistical Mechanics of Lattice Systems}}
\author{}
\date{}

\begin{document}
\maketitle

\noindent
We study the Ising model on $\Z^d$, giving precise definitions and a detailed
analysis of the conditions ensuring the existence or absence of a phase
transition. The plan: \S3.1 defines the model with various boundary conditions;
\S3.2 introduces the thermodynamic limit, the pressure and the magnetization,
computed explicitly for $d=1$ in \S3.3; \S3.4 defines infinite-volume Gibbs
states; \S3.6 develops correlation inequalities (GKS, FKG); \S3.7 analyzes the
phase diagram, proving a phase transition at $h=0$, $d\ge 2$, $\beta$ large
(Peierls), uniqueness at high temperature, and uniqueness at $h\neq0$
(Lee--Yang); \S3.10 collects complements.

% ===========================================================================
\section{Finite-volume Gibbs distributions}
% ===========================================================================
Configurations of the Ising model in a finite volume $\La \fin \Z^d$ with
\emph{free} boundary condition form $\Omega_\La = \{-1,1\}^\La$. The spin at
$i$ is $\sg_i(\omega)=\omega_i$. With $\Eset_\La = \{\,\{i,j\}\subset\La : i\sim
j\,\}$ the nearest-neighbour edges, the energy of $\omega$ is
\[
  \Ham^{\varnothing}_{\La;\beta,h}(\omega)
    = -\beta \sum_{\{i,j\}\in\Eset_\La} \sg_i(\omega)\sg_j(\omega)
      - h \sum_{i\in\La} \sg_i(\omega),
\]
with inverse temperature $\beta\in\Rgeq$ and magnetic field $h\in\R$. The
superscript $\varnothing$ marks the free boundary condition.

\begin{definition}
The Gibbs distribution of the Ising model in $\La$ with free boundary condition
is the probability measure on $\Omega_\La$
\[
  \mu^{\varnothing}_{\La;\beta,h}(\omega)
    = \frac{e^{-\Ham^{\varnothing}_{\La;\beta,h}(\omega)}}
           {Z^{\varnothing}_{\La;\beta,h}},
  \qquad
  Z^{\varnothing}_{\La;\beta,h}
    = \sum_{\omega\in\Omega_\La} e^{-\Ham^{\varnothing}_{\La;\beta,h}(\omega)} .
\]
\end{definition}

\begin{definition}
The Gibbs distribution in the box $V_n$ with \emph{periodic} boundary condition
is defined analogously, with the Hamiltonian summed over the edges of the torus
(nearest-neighbour pairs taken modulo the side length).
\end{definition}

\begin{definition}
The Gibbs distribution of the Ising model in $\La$ with boundary condition
$\eta\in\Omega$ uses the Hamiltonian
\[
  \Ham^{\eta}_{\La;\beta,h}(\omega)
    = -\beta \!\!\sum_{\substack{\{i,j\}\cap\La\neq\varnothing}}\!\!
        \sg_i \sg_j - h\sum_{i\in\La}\sg_i,
\]
where spins outside $\La$ are frozen to $\eta$. The special cases $\eta\equiv+1$
and $\eta\equiv-1$ give the $+$ and $-$ boundary conditions.
\end{definition}

% ===========================================================================
\section{Thermodynamic limit, pressure and magnetization}
% ===========================================================================

\subsection{Convergence of subsets}
We write $\La\incr\Z^d$ (van Hove convergence) for an increasing sequence of
finite volumes invading $\Z^d$ with vanishing relative boundary. The boxes
$B(n)=\{-n,\dots,n\}^d$ are the canonical example.

\subsection{Pressure}
\begin{definition}
The (finite-volume) pressure with boundary condition of type $\bdy$ is
\[
  \psi^{\bdy}_\La(\beta,h)
    = \frac{1}{|\La|}\log Z^{\bdy}_{\La;\beta,h}.
\]
\end{definition}

\begin{lemma}
For each type of boundary condition $\bdy$, the map
$(\beta,h)\mapsto\psi^{\bdy}_\La(\beta,h)$ is convex.
\end{lemma}

\begin{theorem}
In the thermodynamic limit, the pressure
\[
  \psi(\beta,h) = \lim_{\La\incr\Z^d} \psi^{\bdy}_\La(\beta,h)
\]
is well defined, independent of the sequence $\La\incr\Z^d$ and of the type of
boundary condition. Moreover $\psi$ is convex on $\Rgeq\times\R$ and even in $h$.
\end{theorem}
\noindent The proof uses subadditivity for the free boundary condition, then
shows boundary conditions and the choice of sequence do not affect the limit.

\subsection{Magnetization}
The average magnetization density in $\La$ is
$m^{\bdy}_\La(\beta,h)=\frac1{|\La|}\sum_{i\in\La}\avg{\sg_i}^{\bdy}_\La
=\partial_h\psi^{\bdy}_\La$. Let $\bc_\beta$ be the (at most countable) set of
$h$ at which $h\mapsto\psi(\beta,h)$ fails to be differentiable.

\begin{corollary}
For all $h\notin\bc_\beta$, the average magnetization density
$m(\beta,h)=\partial_h\psi(\beta,h)$ exists and equals
$\lim_{\La\incr\Z^d} m^{\bdy}_\La(\beta,h)$, independently of the boundary
condition.
\end{corollary}

\subsection{A first definition of phase transition}
\begin{definition}
The pressure $\psi$ exhibits a \emph{first-order phase transition} at
$(\beta,h)$ if $h\mapsto\psi(\beta,h)$ is not differentiable there, i.e.
$\frac{\partial\psi}{\partial h^-}(\beta,h)
   <\frac{\partial\psi}{\partial h^+}(\beta,h)$.
\end{definition}

% ===========================================================================
\section{The one-dimensional Ising model}
% ===========================================================================
\begin{theorem}[$d=1$]
For all $\beta\ge0$ and $h\in\R$, the pressure of the one-dimensional Ising
model is
\[
  \psi(\beta,h)
    = \log\!\Big\{ e^{\beta}\cosh(h)
        + \sqrt{ e^{2\beta}\cosh^2(h) - 2\sinh(2\beta)}\,\Big\}. \tag{3.10}
\]
In particular $h\mapsto\psi(\beta,h)$ is real-analytic for every $\beta\ge0$,
so $\bc_\beta=\varnothing$: there is \emph{no} phase transition in $d=1$.
\end{theorem}
\noindent The proof uses the transfer-matrix method: $Z$ is a trace of a
product of $2\times2$ matrices, and $\psi$ is the log of the largest eigenvalue.

\begin{theorem}[$d=1$]
Let $0<\beta<\infty$ and $\La_n\incr\Z$ with arbitrary boundary condition. For
all $\epsilon>0$ there is $c=c(\beta,\epsilon)>0$ such that, for $n$ large,
\[
  \mu^{\bdy}_{\La_n;\beta,0}\big( m_{\La_n}\notin(-\epsilon,\epsilon)\big)
    \le e^{-c|\La_n|}. \tag{3.12}
\]
\end{theorem}
\noindent (A large-deviation/concentration statement: at $h=0$ the empirical
magnetization concentrates at $0$.)

% ===========================================================================
\section{Infinite-volume Gibbs states}
% ===========================================================================
\begin{definition}
A function $f:\Omega\to\R$ is \emph{local} if it depends on finitely many spins,
i.e. there is $\Delta\fin\Z^d$ with $f(\omega)=f(\omega')$ whenever
$\omega|_\Delta=\omega'|_\Delta$.
\end{definition}

\setcounter{theorem}{12}% book numbering gap (3.12 not transcribed)
\begin{definition}
An (infinite-volume) \emph{state} is a normalized, positive linear functional
$\avg{\cdot}$ on local functions.
\end{definition}

\begin{definition}
Given $\La_n\uparrow\Z^d$ and boundary conditions $(\bdy_n)$, the associated
infinite-volume Gibbs state is the limit
$\avg{f}=\lim_n \avg{f}^{\bdy_n}_{\La_n;\beta,h}$ for all local $f$, when it
exists.
\end{definition}

\setcounter{theorem}{15}% book numbering gap (3.15 not transcribed)
\begin{definition}
A state $\avg{\cdot}$ is \emph{translation invariant} if
$\avg{f\circ\theta_j}=\avg{f}$ for every local $f$ and every $j\in\Z^d$.
\end{definition}

\begin{theorem}
Let $\beta\ge0$, $h\in\R$. Along any sequence $\La_n\uparrow\Z^d$, the
finite-volume Gibbs distributions with $+$ (resp.\ $-$) boundary condition
converge to infinite-volume Gibbs states
\[
  \meanp{\cdot}_{\beta,h}=\lim_{n}\avg{\cdot}^{+}_{\La_n;\beta,h},
  \qquad
  \meanm{\cdot}_{\beta,h}=\lim_{n}\avg{\cdot}^{-}_{\La_n;\beta,h},
\]
independently of the sequence. These are translation-invariant Gibbs states,
extremal in the sense that $\meanp{\cdot}\ge\avg{\cdot}\ge\meanm{\cdot}$ on
nondecreasing local functions.
\end{theorem}
\noindent The proof relies on monotonicity in the volume and boundary condition
(FKG / Lemma 3.22--3.23) together with the polynomial expansion Lemma 3.19.

% ===========================================================================
\section{Two families of local functions}
% ===========================================================================
\setcounter{theorem}{18}% book numbering gap (3.18 not transcribed)
\begin{lemma}
Every local $f$ admits a finite expansion
$f=\sum_{A\subset\mathrm{supp}(f)} \hat f_A\,\sg_A$ in the monomials
$\sg_A=\prod_{i\in A}\sg_i$, with real coefficients $\hat f_A$. Consequently
agreement of all correlation functions $\avg{\sg_A}$ determines a state.
\end{lemma}
\noindent The two basic monotone families are $\sg_A$ and the occupation
variables $n_i=\tfrac{1+\sg_i}{2}$, $n_A=\prod_{i\in A} n_i$.

% ===========================================================================
\section{Correlation inequalities}
% ===========================================================================

\subsection{The GKS inequalities}
\begin{theorem}[GKS inequalities]
Let $J\ge0$, $h\ge0$ and $\La\fin\Z^d$. Then for all $A,B\subset\La$,
\[
  \meanp{\sg_A}_{\La;J,h}\ge0,\qquad
  \meanp{\sg_A\sg_B}_{\La;J,h}
     \ge \meanp{\sg_A}_{\La;J,h}\meanp{\sg_B}_{\La;J,h}. \tag{3.21--3.22}
\]
The inequalities also hold for free and periodic boundary conditions.
\end{theorem}

\subsection{The FKG inequality}
\begin{theorem}[FKG inequality]
Let $J=(J_{ij})\ge0$, $h=(h_i)\in\R^{\Z^d}$ arbitrary, $\La\fin\Z^d$, and
$\bdy$ any boundary condition. Then for any pair of nondecreasing functions
$f,g$,
\[
  \avg{fg}^{\bdy}_{\La;J,h}
    \ge \avg{f}^{\bdy}_{\La;J,h}\,\avg{g}^{\bdy}_{\La;J,h}. \tag{3.24}
\]
(The inequality reverses if $g$ is nonincreasing.)
\end{theorem}

\subsection{Consequences (monotonicity)}
\begin{lemma}
For nondecreasing $f$ and $\La_1\subset\La_2\fin\Z^d$,
$\avg{f}^{+}_{\La_1}\ge\avg{f}^{+}_{\La_2}$ and
$\avg{f}^{-}_{\La_1}\le\avg{f}^{-}_{\La_2}$: increasing the volume under $+$
b.c.\ decreases the expectation of any nondecreasing $f$ (and conversely under
$-$).
\end{lemma}

\begin{lemma}
For any nondecreasing $f$, $\beta\ge0$ and arbitrary boundary condition $\eta$,
$\meanm{f}_{\La;\beta,h}\le\avg{f}^{\eta}_{\La;\beta,h}\le\meanp{f}_{\La;\beta,h}$.
The $+$ and $-$ states dominate all others. These monotonicity facts are what
make the limits in Theorem 3.17 exist.
\end{lemma}

% ===========================================================================
\section{Phase diagram}
% ===========================================================================
\begin{definition}
If at least two distinct Gibbs states can be constructed for a pair
$(\beta,h)$, one says there is \emph{non-uniqueness} (a phase transition) at
$(\beta,h)$.
\end{definition}

\begin{theorem}[main result]\leavevmode
\begin{enumerate}
  \item In any $d\ge1$, when $h\neq0$ there is a \emph{unique} Gibbs state for
        all $\beta\in\Rgeq$.
  \item In $d=1$, there is a unique Gibbs state at every $(\beta,h)$.
  \item When $h=0$ and $d\ge2$, there exists
        $\beta_c=\beta_c(d)\in(0,\infty)$ such that:
        \begin{itemize}
          \item for $\beta<\beta_c$ the Gibbs state at $(\beta,0)$ is unique;
          \item for $\beta>\beta_c$ a first-order phase transition occurs,
                $\meanp{\cdot}_{\beta,0}\neq\meanm{\cdot}_{\beta,0}$
                (spontaneous magnetization $m^*(\beta)>0$).
        \end{itemize}
\end{enumerate}
\end{theorem}
\noindent The proof is spread over the next sections: item 3 (existence of a
transition) by the Peierls argument and high-temperature uniqueness; items 1--2
by the Lee--Yang analyticity argument.

\subsection{Two criteria for (non)-uniqueness}
\setcounter{theorem}{27}% book numbering gap (3.26--3.27 not transcribed)
\begin{theorem}
For $(\beta,h)\in\Rgeq\times\R$ the following are equivalent:
\begin{enumerate}
  \item there is a unique Gibbs state at $(\beta,h)$;
  \item $\meanp{\cdot}_{\beta,h}=\meanm{\cdot}_{\beta,h}$;
  \item $\meanp{\sg_0}_{\beta,h}=\meanm{\sg_0}_{\beta,h}$.
\end{enumerate}
Thus uniqueness is detected by a single one-point function.
\end{theorem}

\begin{proposition}
For any $\La\incr\Z^d$ the limits $m^{+}(\beta,h)=\meanp{\sg_0}_{\beta,h}$ and
$m^{-}(\beta,h)=\meanm{\sg_0}_{\beta,h}$ exist; $m^{+}$ is right-continuous,
$m^{-}$ left-continuous in $h$, and $m^{-}=-m^{+}$ by spin-flip symmetry.
\end{proposition}

\setcounter{theorem}{30}% book numbering gap (3.30 not transcribed)
\begin{lemma}
For all $\beta\ge0$, $h\mapsto\meanp{\sg_0}_{\beta,h}$ is nondecreasing and
right-continuous (analogous statements for $\meanm{\sg_0}$). This yields the
one-sided continuity of $m^{\pm}$.
\end{lemma}

\begin{definition}
The \emph{critical inverse temperature} is
\[
  \beta_c(d) = \inf\{\beta\ge0 : m^*(\beta)>0\}
             = \sup\{\beta\ge0 : m^*(\beta)=0\},
\]
where $m^*(\beta)=\meanp{\sg_0}_{\beta,0}$ is the spontaneous magnetization. At
$h=0$, uniqueness is equivalent to $m^*(\beta)=0$.
\end{definition}

\setcounter{theorem}{33}% book numbering gap (3.33 not transcribed)
\begin{theorem}
For all $\beta\ge0$, $h\in\R$,
\[
  \frac{\partial\psi}{\partial h^+}(\beta,h)=m^{+}(\beta,h),\qquad
  \frac{\partial\psi}{\partial h^-}(\beta,h)=m^{-}(\beta,h).
\]
Hence $h\mapsto\psi(\beta,h)$ is differentiable at $h$ \emph{iff} there is a
unique Gibbs state at $(\beta,h)$: the analytic notion of first-order transition
(Definition 3.8) coincides with the probabilistic notion of non-uniqueness.
\end{theorem}

\subsection{Spontaneous symmetry breaking at low temperature (Peierls)}
For $h=0$, $d=2$, in the box $B(n)$ with $+$ boundary condition, each
configuration is described by a family of \emph{contours} (closed dual circuits
separating $+$ and $-$ regions).

\setcounter{theorem}{36}% book numbering gap (3.35--3.36 not transcribed)
\begin{lemma}[Peierls' estimate]
For all $\beta>0$ and any contour $\gamma^*$,
\[
  \mu^{+}_{B(n);\beta,0}(\Gamma\ni\gamma^*) \le e^{-2\beta|\gamma^*|}. \tag{3.34}
\]
The probability of a fixed contour decays exponentially in its length: the
ground state $\eta^+$ is \emph{stable} at large $\beta$.
\end{lemma}
\noindent Summing over contours surrounding the origin (and controlling their
number, Lemma 3.38), one gets $\meanp{\sg_0}_{\beta,0}>0$ for $\beta$ large
enough, hence $\beta_c(d)<\infty$ for $d\ge2$, with the explicit bound
$\beta_c(2)<0.88$ (Exercise 3.21).

\begin{lemma}
Combinatorial bound on the number of contours of a given length through a fixed
edge, obtained by an Eulerian traversal of a connected graph: a connected graph
with $N$ edges admits a closed walk using each edge at most twice.
\end{lemma}

\subsection{Uniqueness at high temperature}
A high-temperature (small-$\beta$) expansion represents the two-point function
$\avg{\sg_i\sg_j}^{+}_{\beta,0}$ as a sum over collections of edges/paths. For
$\beta$ small the geometric weight is summable, giving exponential decay of
correlations and therefore $m^*(\beta)=0$, i.e. uniqueness for $\beta<\beta_c$,
and $\beta_c(d)>0$.

\subsection{Uniqueness in nonzero magnetic field}
\setcounter{theorem}{39}% book numbering gap (3.39 not transcribed)
\begin{theorem}
Let $\beta>0$. As a function of $h$, the thermodynamic-limit pressure
$\psi=\psi(h)$ extends from $\{h>0\}$ (resp.\ $\{h<0\}$) to an analytic function
on the right half-plane $H^+$ (resp.\ $H^-$); on $H^\pm$ it can be computed via
the free-boundary thermodynamic limit.
\end{theorem}
\noindent By Theorem 3.34 analyticity in $h$ forces a unique Gibbs state, so
Theorem 3.40 yields item 1 of Theorem 3.25 (uniqueness for all $h\neq0$).

\setcounter{theorem}{41}% book numbering gap (3.41 not transcribed)
\begin{theorem}[Lee--Yang]
Let $\beta\ge0$ and let $D\subset\C$ be open, simply connected, with $D\cap\R$
an interval. If, for every finite volume $\La\fin\Z^d$, the partition function
satisfies $Z^{\varnothing}_{\La;\beta,h}\neq0$ for all $h\in D$, then the
pressure $h\mapsto\psi(h)$ admits an analytic continuation to $D$.
\end{theorem}

\begin{theorem}[Lee--Yang Circle Theorem]
The hypothesis of Theorem 3.42 holds with $D=H^+$ and $D=H^-$. Equivalently, in
the variable $z=e^{-2h}$ all zeros of the partition function lie on the unit
circle $|z|=1$ (i.e.\ at purely imaginary $h$). This is what supplies the
no-zero condition needed for Theorem 3.40.
\end{theorem}

\subsection{Summary of what has been proved}
Combining: a first-order phase transition at $(\beta,0)$ for $\beta>\beta_c(d)$,
$d\ge2$ (Peierls); uniqueness for $\beta<\beta_c(d)$ (high-temperature
expansion); and uniqueness for all $h\neq0$ and all $d$ (Lee--Yang). Together
these establish the full phase diagram of Theorem 3.25.

% ===========================================================================
\section{Proofs of the correlation inequalities}
% ===========================================================================
\S3.8.1 proves the GKS inequalities (Theorem 3.20) by expanding products of
$\sg_A$ and using that all coefficients of $e^{\beta\sg_i\sg_j}$,
$e^{h\sg_i}$ are nonnegative when $\beta,h\ge0$. \S3.8.2 proves the FKG
inequality (Theorem 3.21) from the FKG lattice condition (Holley's inequality /
the Harris--FKG correlation inequality for monotone measures).

% ===========================================================================
\section{Bibliographical references}
% ===========================================================================
Pointers to Ising, Peierls, Griffiths--Kelly--Sherman, Fortuin--Kasteleyn--
Ginibre, Lee--Yang, and the historical literature on the model.

% ===========================================================================
\section{Complements}
% ===========================================================================
A series of more advanced or specialized topics, treated less formally:
\begin{itemize}
  \item[3.10.1] \textbf{Kramers--Wannier duality} — the self-duality of the 2D
        model relating high- and low-temperature regimes and locating the
        self-dual point $\beta_{\mathrm{sd}}$, conjecturally $=\beta_c(2)$.
  \item[3.10.2] \textbf{Mean-field bounds} — comparison inequalities bounding
        $m^*$ above/below by mean-field quantities.
  \item[3.10.3] \textbf{An alternative proof of the FKG inequality.}
  \item[3.10.4] \textbf{Transfer matrix and Markov chains} — the operator
        viewpoint and its probabilistic interpretation.
  \item[3.10.5] \textbf{The Ising antiferromagnet} ($J<0$) and the role of
        bipartiteness.
  \item[3.10.6] \textbf{Random-cluster and random-current representations} —
        the FK/current geometric reformulations, percolation of clusters, and
        the switching lemma.
  \item[3.10.7] \textbf{Non-translation-invariant Gibbs states and interfaces}
        — Dobrushin states in $d\ge3$, rigid interfaces, and their absence in
        $d=2$.
  \item[3.10.8] \textbf{Gibbs states and local behavior in large finite
        systems.}
  \item[3.10.9] \textbf{Absence of analytic continuation of the pressure}
        through the transition point.
  \item[3.10.10] \textbf{Metastable behavior in finite systems.}
  \item[3.10.11] \textbf{Critical phenomena} — critical exponents, scaling, and
        what is/ is not rigorously known.
  \item[3.10.12] \textbf{Exact solution} — Onsager's solution in $d=2$ and the
        explicit $\beta_c(2)=\tfrac12\log(1+\sqrt2)$.
  \item[3.10.13] \textbf{Stochastic dynamics} — Glauber dynamics reversible
        w.r.t.\ the Gibbs measure.
\end{itemize}

\end{document}
